% RESEARCH_DRAFT_SHA256: 5d7a832406f629b7b8d48e403be1a9bd41f9b80947971ce8357ada1a0484e0b3
% Layout and source-reference wrappers expanded for the web renderer.
\documentclass[11pt]{article}
\usepackage[margin=25mm]{geometry}
\usepackage{amsthm}
\usepackage{fontspec,unicode-math}
\setmainfont{Latin Modern Roman}
\setmathfont{Latin Modern Math}
\usepackage{amsmath,amssymb,mathtools,xurl,hyperref}
\hypersetup{colorlinks=true,urlcolor=blue,linkcolor=blue}
\setcounter{secnumdepth}{0}
\setlength{\parindent}{0pt}
\setlength{\parskip}{5pt}
\setlength{\emergencystretch}{3em}
\newcommand{\sref}[2]{\hyperlink{source:#1:#2}{[S#2]}}
\newcommand{\cataloguescope}[1]{\paragraph{Source record.}#1}
\newcommand{\floor}[1]{\left\lfloor#1\right\rfloor}
\newtheorem{proposition}{Proposition}
\newtheorem{lemma}[proposition]{Lemma}
\begin{document}
\paragraph{Research manuscript.} OpenAI \texttt{gpt-6-astra}, reasoning \texttt{ultra}; 11 October 2026.
\paragraph{Preserved source body.} The complete catalogue mathematical text and sources follow with layout and source links expanded, including source qualifications and its historical preparation notes. New research begins at the section headed Research attempt.

% UnsolvedMath ID 20002119; source catalogue code AIM-INFRASTRUCTURE-0019
\section{A trace-validity obstruction and engagement claim certificate}\label{20002119}
{\small UnsolvedMath ID 20002119 · Computer Science}
\subsection{Definitions and mathematical statement}

\paragraph{Shared notation.}Standard number systems are $\mathbb N,\mathbb Z,\mathbb Q,\mathbb R,\mathbb C$. Specialized notation and index ranges follow the cited source and the definitions below.


Educational engagement concerns participation in learning activities, attention to mathematical ideas, and motivation to persist. These are different constructs; a platform log records actions, while an engagement measure operationalizes a selected construct for a population and time period. Enhancing engagement means changing a named measure under a specified use of the technology.
\paragraph{Record scope and qualifications.}
The catalogue title is retained as metadata. Named certificates, contracts, bounds or frameworks in that title are not asserted as results of the original source extract.
No technology, engagement instrument, population, intervention, or comparison condition is specified. The statement is an empirical design question rather than a quantified mathematical assertion.

Catalogue screening limitation: Educational assessment, adoption, or study-design discussion; the supplied question does not state a mathematical target with defined assumptions and outcome.

The source snapshot labels this record partially solved; this is catalogue metadata rather than independent verification of a mathematical result or the current completion of a project.
\paragraph{Statement recorded in the supplied source.}
20. Understand what type of engagement the technology enhances/supports \href{http://aimath.org/pastworkshops/oerassessmentproblems.pdf}{[S2]}
\paragraph{Source status.}
The supplied catalogue labels this source record partially solved; that metadata was checked against the supplied snapshot on 2026-10-11. The label is preserved for provenance and is not an independent verification of present open status or mathematical correctness.
\subsection{Short English statement}
Determine which forms of learner engagement an educational technology supports.
\subsection{Sources}
\begin{itemize}
\item[S1] Ulam AI and the UnsolvedMath Contributors, UnsolvedMath: A Curated Collection of Open Mathematics Problems, record 20002119 (AIM-INFRASTRUCTURE-0019). Catalogue curation and metadata: CC BY 4.0. The source extract is retained; definitions, English summaries and qualifications are editorial additions.
\url{https://huggingface.co/datasets/ulamai/UnsolvedMath}.
\item[S2] AIM Workshop Problem Lists, source linked by catalogue record AIM-INFRASTRUCTURE-0019.
\url{http://aimath.org/pastworkshops/oerassessmentproblems.pdf}.
\end{itemize}
\paragraph{Preparation and review.}
The original source and UnsolvedMath contributors retain attribution for the source extract. Codex (GPT-6; exact serving revision not exposed) assisted with definitions, English summaries, scope notes and formatting. No new human mathematical authorship or independently verified solution is asserted. Permanent source ID 20002119 and code AIM-INFRASTRUCTURE-0019 checked against the complete supplied registry and current repository. Distinct upstream records, including repeated questions and broad or damaged extracts, are retained for complete coverage of the source CS category. A repeated record is not a new independent conjecture. Complete mathematical-equivalence and current-literature review has not been performed. Preparation and final review follow contribution prompts C and D, with source gaps disclosed. Structural and rendering checks do not establish a new theorem or open status.
\label{end:20002119}

\section{Research attempt}
\subsection{Literature, source context and scope}
The primary AIM notes describe a brainstorming session. The original excerpt and all catalogue qualifications remain above. They do not specify a mathematical theorem or a complete empirical study. The formal model below is an explicitly limited research route introduced for this attempt, not a reconstruction of a missing original conjecture. The supplied catalogue's earlier research summary was inspected as context; its named framework is not adopted as an established solution. The elementary result below is not claimed novel, and no educational effectiveness data were collected.

\subsection{Attempt: naming an engagement type from a latent response model}
I tested whether a fitted hidden class of response traces identifies the substantive label ``engaged.'' In each intervention arm $z$, let a binary class $H$ have prevalence $p_z$. Conditional on $H=1$, two binary trace indicators are independent with success probabilities $(a_1,a_2)$; conditional on $H=0$, use $(b_1,b_2)$. The component parameters are shared across arms. This is a limited measurement model, not a definition of cognitive engagement.
\begin{proposition}
The observed trace laws in every arm are unchanged by
\[
(p_z,a,b)\longmapsto(1-p_z,b,a).
\]
Consequently the sign of the change in the prevalence of the class named $H=1$ is not identified without a substantive labeling restriction.
\end{proposition}
\begin{proof}
For a trace vector $y\in\{0,1\}^2$, write
\[
A(y)=\prod_j a_j^{y_j}(1-a_j)^{1-y_j},\qquad
B(y)=\prod_j b_j^{y_j}(1-b_j)^{1-y_j}.
\]
Its observable probability is $p_zA(y)+(1-p_z)B(y)$. The displayed transformation gives $(1-p_z)B(y)+p_zA(y)$, the same expression for every $z,y$. But the transformed prevalence contrast is $(1-p_1)-(1-p_0)=-(p_1-p_0)$. Therefore no statistic of these observable laws can distinguish the two signs while assigning the class label differently.
\end{proof}
\subsection{Verification and remaining gaps}
For $a=(4/5,7/10)$, $b=(1/5,3/10)$ and $(p_0,p_1)=(1/5,4/5)$, all four trace probabilities in each arm were checked under both parameterizations. The putative class-prevalence change is $3/5$ or $-3/5$. This is the familiar label ambiguity of mixtures. Defining one class by higher trace rates would impose a mathematical label, but it would not validate the name cognitive engagement. External anchors or a construct interpretation are needed. No actual technology, engagement instrument or intervention effect is established.

\section{References and source attribution}
American Institute of Mathematics, workshop discussion notes on open educational resources and assessment, \url{https://aimath.org/pastworkshops/oerassessmentproblems.pdf}. The 16-page primary PDF was opened on 11 October 2026 and its numbered brainstorming item was checked against the supplied record. The source supports the provenance and agenda, not the additional modeling assumptions or an empirical educational conclusion. Catalogue and contributor attribution are retained in the reproduced Sources section. The worked mathematics is proved here without reliance on an unverified external theorem.


\paragraph{Integration review.} Codex (GPT-6; exact serving revision not exposed) reviewed the written arguments and expanded layout and source-reference wrappers. Repository metadata and the production web renderer were checked using Python 3.13 and JavaScript with MathJax 3.2.2. This model-assisted review is not independent human proof verification; no PDF compilation or formal proof checking is asserted.

\section{Completion Estimate}
\noindent\textbf{COMPLETION: 0\%}
No substantiated new progress toward the entire original target is claimed. The worked results and failed route are delimited above; a known result, illustrative model, or auxiliary obstruction is not counted as original-target completion.

\section{Final status and attribution}
\textbf{Attempt status: unresolved.} This describes this research attempt, not a new ruling on the source catalogue's status or on external literature claims.
The mathematics and this self-audit were generated by OpenAI \texttt{gpt-6-astra}, reasoning \texttt{ultra}, on 11 October 2026. Codex (GPT-6; exact serving revision not exposed) orchestrated the batch and contributed formatting/repository review. No human mathematical authorship, independent human proof checking, or formal verification is claimed. The historical preparation paragraphs in the reproduced catalogue body concern that body only.

Contribution prompt A was applied to the complete definition and supplied source record, followed by the adversarial checks of prompt B. Local reproduction files are \texttt{batch008/checks.py}, \texttt{computations.json}, and the per-ID audit JSON. They were not public links at preparation time. Numerical checks are finite and do not prove unrestricted results. No LaTeX engine was found locally; compilation was not performed by this research agent.

\end{document}
